\epigraph{Civilization advances by extending the number of important operations which we can perform without thinking about them}{Alfred North Whitehead, An Introduction to Mathematics}

\lettrine{A}{production grade} optimizing compiler tends to be quite a large program.
For example, in LLVM, the peephole rewriter alone is gigantic, and is larger than the vectorizer!
In this thesis, we consider how we go about verifying such a production-grade optimizing compiler.
Naturally, we cannot verify the entire compiler, and thus, we focus on the peephole rewriter,
which is responsible for a small, local transformation on a finite window of instructions.
Each transformation is small enough to prove correct, but the correctness of the entire
rewriter accumulates into thousands of such proof obligations.

Can we prove them by hand? In principle, yes. In practice, this is deathly tedious, error-prone,
which I discovered when embarking on verifying an SSA-based peephole rewriter for LLVM's bitvector optimizations.
To do so, I first formalized a core calculus for SSA-based compiler intermediate representations,
and proved correct a generic peephole rewriter.
After this, I instantiated the generic rewriter against bitvector optimizations, 
and this was where the real challenge began; The proof obligations were complex, subtle,
and it was incredibly tedious to prove them by hand.
The structural half of manipulating SSA was the easy bit.
This tedium drove me to work on for push-button automation,
since it seemed to me like a large bottleneck in being able to verify a production-grade compiler
was the thousands of transformations such a compiler relies on.

What style of decision procedure should we use?
I bet on bitblasting, simply because I had quite a bit of experience thinking about reductions to bitvectors,
and I had an intuition that bitblasting would be a good fit for the kinds of problems we were trying to solve,
which were fundamentally finite, and moreover, about reasoning about bit-level representations.
Unbeknownst to me at the time, this tradition powers today's fastest SMT solvers.
This was a lucky bet, and paid off in ways I certainly didn't appreciate at the time.
For one, I got to learn about the beautiful mathematics of proof certificates, monotone circuits, proof theoretic lower bounds,
k-induction and rIC3, a huge number of intuitions of the structure of a hard search problem,
and the circle of concerns around P, NP, and the hardness of search.
I also came away with a conviction that the only successful algorithms in SMT are those that reduce to SAT,
since this is the one NP-complete problem for which we possess effective algorithms that scale to millions of variables.
I worked on helping implement \bvdecide{}, Lean's fixed-width bitblaster, 
which performs the full end-to-end chain of transforming a bitvector level problem into a SAT instance, solving it,
and then checking the resulting proof certificate.
This infrastructure formed the bedrock of further work in this thesis.

However, the obligations a compiler hands us are not directly bitblastable,
as a rewrite has to be correct at \emph{every} register width, whereas a bitblaster wants one fixed width.
Before my work, the typical solution was to enumerate all widths, and bitblast each one separately.
In this thesis, we show that this is not necessary, and that we can reduce
a family of such width-quantified problems to reachability in a finite automaton,
which we decide by $k$-induction or regular language emptiness.
Interestingly, $k$-induction eventually bottoms out in SAT, which was a fortuitous coincidence at the time.
Next, by a neat encoding trick (where widths are encoded as bitmasks),
any formula over many symbolic widths reduces to an equisatisfiable single-width one.
This translation improves from an exponential number of solver quries to a \emph{single} query,
and additionally, yields a new sound and complete fragment of the multi-width theory.

Finally, I decided to work on floating point, since it was the next missing chunk for compiler verification.
Toward this, I mechanized the SMT-LIB \qffp{} theory and build verified bitblasting circuits for almost all of its operations,
and tweaked the existing proofs of floating point bitblasting to work for all combinations of exponent and mantissa widths.
However, at the end of this effort, I am no longer certain that bitblasting floating point via the IEEE semantics is the right approach.
Instead, I have been pondering an alternative approach based on exponential rings, which I explain in future work.

Overall, did the bet of building proof automation pay off? I certainly think so.
For example, our width-generic solvers prove more all-width problems than the unverified state of the art,
and our floating-point bitblaster, while slower than Bitwuzla,
offers end-to-end guarantees down to Lean's kernel in exchange.
The price of verification does not bite on peephole-sized problems.
Overall, what I learnt was a flavour of the bitter lesson: the only thing that scales is search,
and that the best verification algorithms are ones that reduce to SAT.
While I ache with the loss of the beauty of interpolation,
I grudgingly admit that the best-known algorithms for SMT typically reduce to SAT.
Moreover, I learnt that mechanizing such algorithms is within reach of today's proof assistants,
and foundational proof and push-button automation need not be at odds.
I hope that these reductions lay the groundwork for verification efforts of far greater ambition.
